\section{Exchange retuning and the effective probe}

\subsection{Scalar and quadrupolar local factors}

At a form-factor zero, virtual excitations outside the local ground
multiplet determine the leading scattering amplitude. Let $P=\prod_n P_n$
and $R=(1-P)/(H_0-NE_5)$, where $E_5$ is the isolated ground-multiplet
energy. At $A(q,\alpha_*)=0$, the effective probe becomes
\begin{align}
 (O_q)_{\mathrm{eff}}&=\beta K_q+O(\beta^2),\notag\\
 K_q&=-P O_q R V P-P V R O_q P.
 \label{eq:rank-probe}
\end{align}
The two terms describe a virtual excitation created by the probe or by
the intercell exchange. They are evaluated separately for non-Hermitian $O_q$.

A link can contribute only through an intermediate excitation on the cell
acted on by the probe: the local probe cannot de-excite both cells of a link.
For a Hermitian component of $B^z$, the relevant local tensor is
\begin{equation}
 M^{z\mu}=P_n(B^z R_n S_{\mathrm{end}}^\mu
                   +S_{\mathrm{end}}^\mu R_n B^z)P_n.
 \label{eq:rank-local-tensor}
\end{equation}
It is Hermitian and even under time reversal. Coupling two vectors permits
ranks zero, one, and two, but a Hermitian rank-one tensor within a single
spin irrep is odd under time reversal. Only the scalar and quadrupolar parts
remain. The real and imaginary Hermitian components of $B(q)$ can be treated
separately and recombined.

Define
\begin{equation}
 Q_n^{z\mu}=\tfrac12\{T_n^z,T_n^\mu\}
        -\tfrac13 s(s+1)\delta_{z\mu}I.
\end{equation}
On one orientation of a link the correction consequently has the form
\begin{equation}
 \delta B_n^z=a\,T_{n+1}^z
       +b\sum_\mu Q_n^{z\mu}T_{n+1}^\mu,
 \label{eq:rank-link}
\end{equation}
with a corresponding term for the opposite endpoint.
The second term couples a local quadrupole to a neighboring spin,
forming a composite vector under simultaneous spin rotations.

\subsection{Cancellation for a local doublet}

In a single spin-$1/2$ irrep, $Q_n^{z\mu}=0$ identically. Equivalently every
Hermitian local matrix is a linear combination of $I$ and $\bm T$; time
reversal leaves only $I$ in Eq.~\eqref{eq:rank-local-tensor}.
The entire first-order correction is thus a sum of neighboring one-spin
operators. On a uniform periodic chain at a commensurate momentum, translation
reduces this sum to $K_q=\gamma(q)D_Q$. Cell-centered reflection and real
exchange give $a_R=a_L^*$ for the two link orientations. Consequently
$\gamma=a_L e^{-iQ}+a_R e^{iQ}=2\operatorname{Re}(a_L e^{-iQ})$
is real in the same phase convention as $A$.

Suppose the local zero is simple, with
$A'=\partial_\alpha A(q,\alpha_*)\ne0$. A physical exchange retuning
$\alpha=\alpha_*+\kappa\beta$ gives
\begin{equation}
 (O_q)_{\mathrm{eff}}
   =\beta[\gamma+\kappa A']D_Q+O(\beta^2).
\end{equation}
Choosing $\kappa=-\gamma/A'$ cancels the first-order operator throughout
the retained Hilbert space. Every low-band transition therefore satisfies
\begin{equation}
 \Wlow=O(\beta^4)\quad\text{or smaller}.
 \label{eq:rank-doublet}
\end{equation}
This result applies to a single isolated doublet with SU(2) and
time-reversal symmetry, bilinear links, and the stated periodic and
reflection symmetries. Higher-order retuning can further reduce the
remaining weight; open endpoints are treated in
Appendix~\ref{app:ground-controls}.

\subsection{Residual for a quartet}

For $s=3/2$, the composite vector supplies an additional scattering
amplitude. To determine whether it survives in the ground-state response,
define $d=D_Q|G\rangle$ and
$k=K_q|G\rangle$, with the entire ground subspace projected out.
For degenerate ground states these vectors denote a direct sum, with norm
averaged over the ensemble used in Eq.~\eqref{eq:rank-weight}.
The leading weight after a real retuning is
\begin{equation}
 \Wlow=\beta^2\|k+\kappa A'd\|^2+O(\beta^3).
\end{equation}
Elementary least squares gives
\begin{align}
 \kappa_{\mathrm{opt}}
   &=-\frac{\operatorname{Re}\langle d,k\rangle}{A'\|d\|^2},\\
 \Cmin&=\|k\|^2-
       \frac{[\operatorname{Re}\langle d,k\rangle]^2}{\|d\|^2}.
 \label{eq:rank-minimum}
\end{align}
The minimum is the part of the source orthogonal to the accessible
dipole correction. It depends on the many-body state as well as on the
local tensor structure. The overlaps in the examples below are real to
numerical precision.

At $q=\pi$ the quartet branch provides a simple analytic example:
\begin{align}
 \alpha_*&=-\frac{2+\sqrt7}{6},\notag\\
 c_1=c_5&=\frac{25+4\sqrt7}{108},\quad
 c_2=c_4=\frac14,\notag\\
 c_3&=\frac{1-2\sqrt7}{27}.
 \label{eq:rank-exact-root}
\end{align}
These coefficients obey $c_1-c_2+c_3-c_4+c_5=0$.
Here $\Dloc=0.16622155$, $c_1^2=0.10855198$, and
$A'=0.10026088$. The local resolvent gives $a=-0.12123577$ and
$b=0.12361176$, with exact expressions and a proof that $b>0$ in
Appendix~\ref{app:exact-local-check}. Thus this root realizes a nonzero
composite contribution within the microscopic pentamer model.

\subsection{A bond-spin interpretation}

The composite vector also has a direct interpretation in terms of neighboring
total-spin channels. For two spin-$s$ sites, put
$h_{ij}=\bm T_i\cdot\bm T_j$, $D_{ij}=T_i^z-T_j^z$, and
$X_{ij}=\sum_\mu(Q_i^{z\mu}T_j^\mu-T_i^\mu Q_j^{z\mu})$.
Expanding the quadrupoles gives the exact identity
\begin{equation}
 X_{ij}=\tfrac12\{h_{ij},D_{ij}\}+\tfrac13s(s+1)D_{ij}.
 \label{eq:bond-identity}
\end{equation}
Within bond total spin $J$, $h_{ij}$ has eigenvalue
$[J(J+1)-2s(s+1)]/2$. The exchange-odd vector $D_{ij}$ connects only
adjacent bond-spin sectors. If $\Pi_J$ denotes their projector, then
\begin{equation}
 \Pi_J X_{ij}\Pi_{J-1}=
 \left[\frac{J^2}{2}-\frac{2s(s+1)}{3}\right]
 \Pi_J D_{ij}\Pi_{J-1}.
 \label{eq:bond-channel}
\end{equation}
For a doublet there is only $J=1\leftrightarrow0$, and this coefficient
vanishes. For a quartet, the $1\leftrightarrow0$, $2\leftrightarrow1$,
and $3\leftrightarrow2$ coefficients are $-2,-1/2,2$.
The link probe $-aD_{ij}+bX_{ij}$ therefore renormalizes these channels
differently, whereas exchange retuning adds the same multiple of $D_{ij}$.
One adjustment cannot cancel all three link-operator coefficients.
The ground state can nevertheless sample a restricted combination of
these channels and remain dark. Thermal occupation brings additional
initial states into play, making simultaneous cancellation a property
of the full low-energy dynamics.
